\SourceLeaf{158}
\begin{center}
\footnotesize\setlength{\tabcolsep}{3pt}
\begin{tabular}{l r|r}
&Zinsfuss \(3\tfrac{1}{2}\) Proc.&Zinsfuss 4 Proc.\\
\shortstack[l]{Reine Verbindlichkeit der Kasse gegen\\42 verheirathete Mitglieder}&50008 Thl.&44482 Thl.\\
\shortstack[l]{Danach verhältnissmässig gegen\\9 unverheirathete Mitglieder . . . .}&10716 ,, &9532 ,, \\[5pt]
\shortstack[l]{Totalwerth der zweiten Rubrik\\ . . . . . . . . . . . . . . . . . . . . . . . .}&60724 ,, &54014 ,,
\end{tabular}
\end{center}
\noindent\hspace*{\parindent}Dass einer solchen Evaluirung nicht ganz dieselbe Zuverlässigkeit beigelegt werden kann, wie der Berechnung der Verbindlichkeiten der ersten Rubrik, habe ich schon in meinem mehrerwähnten Gutachten bevorwortet. Es finden nemlich bei unsrer Witwenkasse mehrere Umstände Statt, deren Wirkung einer Vorausberechnung gar nicht fähig ist, von welcher aber auch, in Ermangelung einer solchen Buchführung für die frühere Zeit, wie S. [145 unten] erwähnt ist, nicht einmal eine Schätzung gemacht werden kann. Als den einflussreichsten dieser Umstände bezeichne ich die Wiederverheirathung von Mitgliedern nach dem Tode ihrer Ehefrauen. Das Ableben der Frauen vor den Männern ist nemlich schon ein Theil der in der Rechnung berücksichtigten Chancen, und so ist jede Wiederverheirathung eine ausserhalb der Rechnung liegende neu hinzukommende Belastung des Kassen-Conto. Von der andern Seite kommt diesem Kassen-Conto wie eine nicht veranschlagte Entlastung zu Gute jeder anderweitige Abgang eines verheiratheten Mitgliedes, z. B. durch eine auswärtige Vocation. Auch der Umstand, dass solche verwitwete Mitglieder, die sich nicht wieder verheirathen, doch nach dem Tode der Frauen wenigstens eine Zeitlang die Beiträge fortzuzahlen pflegen, gehört in dieselbe Kategorie, obwohl er von geringer Erheblichkeit ist. Endlich ist vielleicht die Rechnung für die dermalen unverheiratheten Theilnehmer nach dem Durchschnittsresultate für die Verheiratheten, etwas zu hoch. Der jedesmalige Bestand der unverheiratheten Mitglieder wird gewöhnlich so zusammengesetzt sein, dass mit Wahrscheinlichkeit angenommen werden kann, ein oder der andere davon werde sich überall nicht verheirathen: diess wird jedoch, wenigstens zum Theil, wieder dadurch aufgewogen, dass, je später Verheirathung erfolgt, desto mehr mit Wahrscheinlichkeit das künftige Überleben der Frau und ein langer Witwenstand präsumirt werden muss.

Bei der Unmöglichkeit, die Wirkung dieser verschiedenen Umstände im Voraus in Zahlen zu veranschlagen, darf man doch für wahrscheinlich halten, dass sie einander grossentheils compensiren.

\Needspace{6\baselineskip}\pdfbookmark[3]{III. Evaluirung der Verbindlichkeiten gegen künftige Mitglieder}{witwen-II-2-III}
\par\medskip\noindent III.\quad Evaluirung des jetzigen Geldwerths der Verbindlichkeiten der Kasse gegen alle künftig beitretenden Mitglieder.\par\smallskip
Für diese dritte Rubrik bieten die bisherigen Erfahrungen der Kasse ein sehr schätzbares Hülfsmittel dar; ja ohne diese Erfahrungen würde eine Veranschlagung ganz unmöglich sein. Mein Verfahren ist folgendes.

Ich habe für jeden der bis Michaelis 1804 beigetretenen Theilnehmer, bis einschliesslich Nro. 108\footnote{Dass ich gerade soweit und nicht weiter gegangen bin, ist mit Vorbedacht geschehen, es würde mich aber zu weit führen, wenn ich die Gründe hier ausführlich entwickeln wollte.}, sowohl den Werth der Beiträge, als den Werth der von ihren Hinterbliebenen bezogenen Pensionen auf den Zeitpunkt ihres Beitritts reducirt, und nach zweierlei Zinsfuss discontirt, jedoch mit folgenden, für den Gebrauch, der von diesen Rechnungen gemacht werden sollte, nothwendigen Modificationen:

1) Die jährlichen Beiträge und die Pensionen sind in Rechnung gebracht, nicht wie sie wirklich gezahlt sind, sondern nach ihrer jetzigen Höhe, nemlich 10 Rthl. für jene, 250 Rthl. für diese.

2) Die Waisenpensionen, welche bis 1794 nur bis zum vollendeten 12\textsuperscript{ten} Jahre des jüngsten Kindes verabreicht wurden, sind so gerechnet, als ob sie noch 8 Jahre länger gedauert hätten. Durch Nachfor%
